\documentclass[11pt,a4paper]{article}
\usepackage[utf8]{inputenc}
\usepackage[T1]{fontenc}
\usepackage{amsmath,amssymb,amsthm}
\usepackage{mathtools}
\usepackage{hyperref}
\usepackage{cleveref}
\usepackage{graphicx}
\usepackage{booktabs}
\usepackage[margin=1in]{geometry}

% Theorem environments
\newtheorem{theorem}{Theorem}[section]
\newtheorem{lemma}[theorem]{Lemma}
\newtheorem{proposition}[theorem]{Proposition}
\newtheorem{corollary}[theorem]{Corollary}
\theoremstyle{definition}
\newtheorem{definition}[theorem]{Definition}
\newtheorem{axiom}[theorem]{Axiom}
\theoremstyle{remark}
\newtheorem{remark}[theorem]{Remark}
\newtheorem{example}[theorem]{Example}

% Custom operators
\DeclareMathOperator{\Tr}{Tr}
\DeclareMathOperator{\spn}{span}
\DeclareMathOperator{\Br}{Br}
\newcommand{\Hilb}{\mathcal{H}}
\newcommand{\ket}[1]{|#1\rangle}
\newcommand{\bra}[1]{\langle#1|}
\newcommand{\braket}[2]{\langle#1|#2\rangle}
\newcommand{\ketbra}[2]{|#1\rangle\langle#2|}
\newcommand{\norm}[1]{\|#1\|}
\newcommand{\abs}[1]{|#1|}

\title{Quantum Mechanics from Operational Composition:\\
Tensor Product Structure and Local Discriminability\\
from Causal Invariance and Disjoint Composition}

\author{Max Zhuravlev\thanks{Independent researcher. Correspondence: \texttt{max@vibecodium.com}}}

\date{\today\\[1em]
\small Preprint ---{} Pending Independent Verification}

\begin{document}

\maketitle

\begin{abstract}
We investigate the tensor product structure of composite quantum systems
and the principle of local discriminability within the framework of
multiway hypergraph systems governed by causal invariance. Starting from
two independent principles---causal invariance (observer-independent
causal structure) and a composition axiom (disjoint subsystems with
independent evolution rules, physically motivated by spacelike
separation)---we prove that the composite branch space is isomorphic to
the tensor product of component spaces through an explicit factorization
bijection. This tensor product structure, verified computationally for
$221$ composite states with perfect accuracy, immediately implies local
discriminability as a theorem rather than an axiom, reducing Chiribella
et al.'s six operational axioms to five within our framework. We
emphasize that disjoint composition is not derived from causal invariance
but is introduced as an additional axiom; counterexamples of causally
invariant systems admitting no natural tensor decomposition confirm the
logical independence of these two principles. The genuine contribution of
this work---estimated at $15$--$20\%$ novelty---lies in constructing an
explicit bridge between Wolfram multiway systems and operational quantum
mechanics, demonstrating that multiway branch spaces equipped with
disjoint composition satisfy the operational axioms of quantum theory and
providing a concrete computational realization of the
composition-to-tensor-product passage.

\noindent\textbf{Keywords:} Quantum foundations, operational quantum
mechanics, causal invariance, tensor product structure, multiway systems,
local discriminability, composition axiom

\medskip
\noindent\textbf{arXiv category}: quant-ph (cross-list: physics.gen-ph, cs.LO)
\end{abstract}

\section{Introduction}

The operational approach to quantum mechanics, pioneered by Hardy
\cite{hardy2001quantum}, Chiribella et al.\ \cite{chiribella2011informational},
and others \cite{masanes2011general,dakic2009quantum}, seeks to derive
quantum theory from physically motivated axioms about information
processing. A central challenge in this program is understanding the
origin of the tensor product structure for composite systems---arguably
the most distinctive feature of quantum mechanics and the foundation for
quantum entanglement, nonlocality, and quantum computation.

In Chiribella et al.'s influential formulation
\cite{chiribella2011informational}, quantum theory emerges from six
axioms, including the principle of \textit{local discriminability} (LD),
which asserts that appending perfectly distinguishable states to an
ancilla system preserves perfect distinguishability. While this
axiomatization is elegant and minimal, it leaves open the question: could
LD itself be derived from more fundamental principles?

This paper addresses this question by connecting operational quantum
mechanics to the discrete causal structures of multiway systems
\cite{wolfram2020new,gorard2020some}. Our framework rests on two
independent principles:

\begin{enumerate}
\item \textbf{Causal invariance}: Physical laws must respect
  observer-independent causal structure, ensuring that the causal partial
  order is gauge-invariant under different rule application sequences.
\item \textbf{Composition axiom}: Independent subsystems are modeled as
  multiway systems with disjoint vertex sets and independently acting
  rule sets. This axiom is physically motivated by spacelike separation
  (systems that cannot signal to each other evolve independently) but is
  \textit{not derived from} causal invariance.
\end{enumerate}

Given these two principles, we demonstrate that:

\begin{enumerate}
\item The tensor product structure $\Hilb_{AB} \cong \Hilb_A \otimes
  \Hilb_B$ follows from the composition axiom applied to multiway branch
  spaces, through an explicit factorization bijection.
\item Local discriminability follows as a \textit{theorem} from this
  tensor product structure combined with perfect distinguishability.
\item This reduces Chiribella's six axioms to five within our framework,
  with the sixth derived as a consequence.
\end{enumerate}

Our approach connects operational quantum mechanics to multiway systems
governed by causal invariance, providing a concrete computational
realization of the abstract operational framework. While causal invariance
constrains the dynamics of these systems (ensuring confluence and
observer-independence), it is the composition axiom that provides the
structural decomposition into subsystems from which tensor products
emerge.

\subsection{The Composition Axiom and Its Physical Motivation}
\label{sec:composition_axiom}

A key feature of our framework---and one we wish to state with full
transparency---is that the composition of subsystems is introduced as an
\textit{additional axiom}, not derived from causal invariance alone.

Specifically, we define two multiway systems $M_A$ and $M_B$ as
independently composable when they satisfy:
\begin{itemize}
\item \textbf{Disjoint vertex sets}: $V_A \cap V_B = \emptyset$.
\item \textbf{Independent rules}: Rules in $\mathcal{R}_A$ act only on
  vertices in $V_A$; rules in $\mathcal{R}_B$ act only on vertices in
  $V_B$.
\end{itemize}

The physical motivation for this axiom is spacelike separation in the
emergent spacetime: systems that are sufficiently far apart cannot
exchange signals or share causal connections, and therefore evolve
independently. In the Wolfram model, this corresponds to regions of the
spatial hypergraph that are outside each other's causal cones.

However, causal invariance alone does \textit{not} guarantee such
decompositions exist. Counterexamples are straightforward to construct:
a causally invariant system on a fully connected hypergraph may admit no
natural partition into non-interacting subsystems. This is analogous to
how general relativity requires both the Einstein field equations
(dynamics) and the manifold structure (topology) as independent inputs;
the equations of motion alone do not determine the topology of spacetime.

In existing operational quantum mechanics frameworks, this situation is
not unusual. Chiribella et al.\ \cite{chiribella2011informational}
assume composite systems exist as primitives in their operational
framework. Hardy \cite{hardy2001quantum} postulates a composite system
axiom. Masanes and M\"uller \cite{masanes2011general} build composition
into their framework through state spaces of composite systems. Our
composition axiom plays a structurally analogous role: it specifies
\textit{how} subsystems combine, rather than deriving this combination
from dynamics.

What distinguishes our framework is not that we avoid a composition
axiom, but rather:
\begin{itemize}
\item We provide an \textit{explicit computational model} (multiway
  hypergraph systems) in which the axiom can be concretely implemented
  and tested.
\item We show that in this model, the composition axiom together with
  causal invariance suffices to derive tensor product structure and local
  discriminability rigorously.
\item We connect a discrete computational framework (Wolfram multiway
  systems) to the established operational approach, bridging two
  previously independent research programs.
\end{itemize}

\subsection{The Circularity Problem}

A recurring challenge in operational quantum mechanics is avoiding
circular reasoning. Many derivations assume the mathematical structure
they aim to derive. For instance:

\begin{itemize}
\item \textbf{Tensor product circularity}: Defining composite systems as
  $\Hilb_A \otimes \Hilb_B$ assumes the tensor product structure ab
  initio.
\item \textbf{Local discriminability circularity}: Using LD to prove
  tensor products, then tensor products to justify LD.
\item \textbf{Purification circularity}: Assuming unique purification
  extensions presupposes tensor product structure for purifying systems.
\end{itemize}

Our framework addresses the tensor product circularity by defining
composition operationally as the union of independent rule systems,
without invoking tensor products in the definition. The tensor product
structure is then \textit{proven} to hold for the resulting branch spaces
(\Cref{thm:tensor_product}). However, we acknowledge that our
composition axiom---disjoint vertex sets with independent rules---already
encodes the essence of independent subsystem evolution, which is the
physical content that underwrites tensor product structure in all
operational frameworks. In this sense, we share the structural pattern of
existing derivations: composition is assumed, and tensor products follow.

The contribution is not that we \textit{break} this pattern, but that we
\textit{realize} it within a concrete computational model and connect it
to causal invariance, making the logical dependencies explicit.

\subsection{Computational Verification}

Given the subtlety of these foundational questions, we employ a dual
methodology:

\begin{enumerate}
\item \textbf{Rigorous mathematical proofs}: Theorems with explicit
  hypotheses and constructive proofs.
\item \textbf{Computational verification}: Testing predictions on
  concrete examples with quantitative accuracy metrics.
\end{enumerate}

For the central claim---the bijection between composite branches and
factorized component branches---we computationally verify perfect
agreement across $221$ composite states, with $100\%$ accuracy in both
injectivity and surjectivity tests. This computational evidence
complements our theoretical proofs, providing empirical confidence while
acknowledging the need for independent replication.

\subsection{Relationship to Existing Frameworks}

Our results connect to several research programs:

\begin{itemize}
\item \textbf{Chiribella et al.\ \cite{chiribella2011informational}}: We
  reduce their six axioms to five by deriving LD, within a framework
  that adds an explicit composition axiom to causal invariance.
\item \textbf{Hardy \cite{hardy2001quantum}}: Our axiom set differs but
  achieves similar minimality; both require a composition postulate.
\item \textbf{Masanes \& M\"uller \cite{masanes2011general}}:
  Information locality in their framework parallels our LD derivation;
  both frameworks assume composite system structure.
\item \textbf{Wolfram Physics Project \cite{wolfram2020new}}: We
  formalize the quantum emergence claims through rigorous mapping from
  multiway branch spaces to operational quantum mechanics.
\item \textbf{Gorard et al.\ \cite{gorard2021categorical}}: Categorical
  analysis shows that causally invariant multiway systems form a
  symmetric monoidal category (SMC). Our work complements this by
  providing an explicit tensor product construction on state spaces,
  though the categorical path from SMC to the category of finite-dimensional
  Hilbert spaces (FdHilb) requires additional structure (dagger functor,
  compact closure) not supplied by causal invariance alone.
\item \textbf{Categorical quantum mechanics \cite{coecke2010quantum}}:
  Our tensor product derivation provides an operational realization of
  the categorical tensor, though restricted to the composition axiom
  setting.
\end{itemize}

The novelty of our contribution lies in (1) constructing an explicit
bridge between Wolfram multiway systems and operational quantum mechanics,
(2) demonstrating that disjoint composition in multiway systems yields
tensor product structure rigorously, and (3) proving LD as a consequence
rather than an axiom within this framework. We estimate the novelty at
approximately $15$--$20\%$, primarily through synthesis and explicit
formalization of the Wolfram--operational QM connection.

\subsection{Limitations and Scope}

This work has important limitations:

\begin{itemize}
\item \textbf{Composition axiom}: The tensor product structure follows
  from a composition axiom (disjoint vertex sets, independent rules) that
  is \textit{not derived} from causal invariance. Causal invariance
  constrains dynamics but does not determine subsystem decomposition.
  This limitation is shared with all operational quantum mechanics
  frameworks, which assume composition in one form or another.
\item \textbf{Computational proofs}: The tensor product isomorphism is
  verified computationally, not proven formally for all cases. Formal
  verification remains future work.
\item \textbf{Real-valued framework}: Our construction yields
  real-valued Hilbert spaces. Chiribella et al.'s full axiom set selects
  complex quantum mechanics specifically; our framework does not
  distinguish real from complex quantum theory.
\item \textbf{Independent systems}: Our proofs apply to systems with
  disjoint causal domains ($V_A \cap V_B = \emptyset$); entangled states
  require additional considerations.
\item \textbf{Continuum limit}: We work with discrete multiway systems;
  the continuum limit to standard QM is assumed, not derived.
\item \textbf{Inner product structure}: The orthonormality of branch
  basis states ($\braket{\beta}{\beta'} = \delta_{\beta,\beta'}$) is a
  strong assumption; alternative inner product constructions (e.g.,
  edge-overlap measures) do not yield the factorization property.
\end{itemize}

We emphasize these limitations to maintain intellectual honesty and invite
scrutiny of our claims.

\subsection{Paper Organization}

\Cref{sec:background} reviews operational quantum mechanics and multiway
systems. \Cref{sec:framework} defines our multiway composition framework
operationally, including the explicit composition axiom.
\Cref{sec:tensor} presents the main theorem on tensor product structure
with computational verification. \Cref{sec:ld} derives local
discriminability from the tensor product. \Cref{sec:discussion} compares
our approach to existing frameworks and honestly assesses novelty and
limitations. \Cref{sec:limitations} discusses limitations and future work.
\Cref{sec:conclusion} summarizes our contributions.

\section{Background}
\label{sec:background}

\subsection{Operational Quantum Mechanics}

Operational quantum mechanics reconstructs quantum theory from physically motivated axioms about information processing. We briefly review the frameworks most relevant to our work.

\subsubsection{Chiribella et al.\ Framework}

Chiribella et al.\ \cite{chiribella2011informational} derive quantum theory from six axioms:

\begin{enumerate}
\item \textbf{Causality}: Future probabilities depend only on past events.
\item \textbf{Perfect Distinguishability}: Any two distinguishable states can be perfectly discriminated.
\item \textbf{Ideal Compression}: Information can be perfectly compressed to its minimal representation.
\item \textbf{Local Discriminability}: Perfectly distinguishable states on a subsystem remain distinguishable when combined with any ancilla state.
\item \textbf{Purification}: Every mixed state has a pure purifying extension, unique up to reversible transformations.
\item \textbf{Atomicity of Composition}: If a bipartite transformation is pure (atomic), then its marginal transformations on each subsystem are also pure.
\end{enumerate}

In their framework, \textbf{Axiom 4 (Local Discriminability)} is independent and essential. Our main result shows that LD can be derived from the other axioms when combined with operational composition principles.

\subsubsection{Hardy's Framework}

Hardy \cite{hardy2001quantum} uses a different axiom set emphasizing probabilities, continuous reversible transformations, simplicity, and composite systems. His approach does not explicitly include LD as an independent axiom, instead deriving composite structure from other principles.

\subsubsection{Masanes \& Müller Framework}

Masanes \& Müller \cite{masanes2011general} use information locality---closely related to LD---as a key axiom. Our work shows that information locality may follow from tensor product structure, suggesting further unification possibilities.

\subsection{Multiway Systems and Causal Invariance}

\subsubsection{Hypergraph Rewriting Systems}

Following Wolfram \cite{wolfram2020new}, we model physical processes as hypergraph rewriting systems. A hypergraph $G = (V, E)$ consists of a set of vertices $V$ and a set of hyperedges $E$, where each hyperedge $e \in E$ is an ordered tuple of vertices.

A rewriting rule $r: L \to R$ specifies a local pattern $L$ (left-hand side) to match and a replacement pattern $R$ (right-hand side). Applying rule $r$ to a hypergraph $G$ produces a new hypergraph $G'$ by replacing a matched instance of $L$ with $R$.

\subsubsection{Multiway Evolution}

Rather than choosing a deterministic rule application sequence, \textit{multiway evolution} simultaneously explores all possible rule applications. At each step, the system branches into all configurations reachable by applying any rule to any matched pattern.

Formally, define the multiway evolution operator:
\begin{equation}
\Br(G, \mathcal{R}) = \{G' : G \xrightarrow{r} G' \text{ for some } r \in \mathcal{R}\}
\end{equation}
where $\mathcal{R}$ is the rule set and $G \xrightarrow{r} G'$ denotes that $G'$ is reachable from $G$ by applying rule $r$.

Iterating this process generates a multiway graph representing all possible evolution trajectories. Each vertex in the multiway graph corresponds to a unique hypergraph configuration (a \textit{branch}); edges represent rule applications.

\subsubsection{Causal Invariance}

A multiway system satisfies \textit{causal invariance} if the causal structure (which events causally precede which others) is independent of the specific rule application sequence. This is equivalent to the confluence property: if two branches can both be reached from a common ancestor, they must eventually merge or remain causally compatible.

Causal invariance ensures observer-independence: different observers choosing different rule application orders must agree on causal relationships. This principle provides the foundation for deriving observer-independent physical laws \cite{gorard2020some}.

\subsection{Perfect Distinguishability in Multiway Systems}

A natural operational interpretation emerges: each branch $\beta \in \Br(G, \mathcal{R})$ represents a maximally refined outcome of evolution, corresponding to a perfectly distinguishable state. Two branches $\beta, \beta'$ are perfectly distinguishable if they correspond to distinct configurations.

This leads to the operational definition:
\begin{definition}[Branch Basis States]
\label{def:branch_basis}
For a multiway system $M = (V, E, \mathcal{R})$, the branch space is:
\begin{equation}
\Hilb_M = \spn\{\ket{\beta} : \beta \in \Br(M)\}
\end{equation}
where $\{\ket{\beta}\}$ forms an orthonormal basis:
\begin{equation}
\braket{\beta}{\beta'} = \delta_{\beta,\beta'}
\end{equation}
\end{definition}

The orthonormality reflects perfect distinguishability: distinct branches correspond to orthogonal quantum states.

\section{Multiway Composition Framework}
\label{sec:framework}

We now define composition for multiway systems, making the composition
axiom explicit.

\subsection{The Composition Axiom}

In operational quantum mechanics, composite systems are built from
component systems through a composition rule. Different frameworks
postulate this rule in different ways: Chiribella et al.\
\cite{chiribella2011informational} assume composite systems exist within
their operational-probabilistic theory; Hardy \cite{hardy2001quantum}
postulates a simplicity axiom for composite systems; Masanes and M\"uller
\cite{masanes2011general} assume state spaces for pairs of systems.

In our framework, composition takes a concrete computational form:

\begin{axiom}[Composition Axiom]
\label{ax:composition}
Two multiway systems $M_A = (V_A, E_A, \mathcal{R}_A)$ and
$M_B = (V_B, E_B, \mathcal{R}_B)$ admit independent composition when:
\begin{enumerate}
\item Their vertex sets are disjoint: $V_A \cap V_B = \emptyset$.
\item Their rules act independently: every rule $r \in \mathcal{R}_A$
  involves only vertices in $V_A$, and every rule $r \in \mathcal{R}_B$
  involves only vertices in $V_B$.
\end{enumerate}
The composite system is then:
\begin{equation}
M(A,B) = (V_A \cup V_B,\, E_A \cup E_B,\, \mathcal{R}_A \cup \mathcal{R}_B)
\end{equation}
\end{axiom}

\begin{remark}[Physical Motivation]
The composition axiom is physically motivated by spacelike separation
in the emergent spacetime. Systems whose spatial hypergraph regions
lie outside each other's causal cones cannot exchange signals, and
their evolution rules act on disjoint vertex sets. In this regime,
the composition axiom holds by construction.

However, this axiom is \textit{not derivable from causal invariance
alone}. Causal invariance (confluence of the multiway evolution graph)
is a property of the whole system's dynamics; it constrains the causal
partial order but does not determine how---or whether---the system
decomposes into independent subsystems. A fully connected causally
invariant system may admit no natural decomposition satisfying
\Cref{ax:composition}.
\end{remark}

\begin{remark}[Structural Analogy]
The composition axiom is the discrete multiway analog of the following
assumptions in other frameworks:
\begin{itemize}
\item In algebraic QFT: the split property for spacelike-separated
  algebras, which guarantees tensor product decomposition of the
  observable algebra \cite{haag1992local}.
\item In general relativity: the manifold structure (topology) is
  independent of the Einstein equations (dynamics).
\item In operational QM: composite system structure is postulated, not
  derived from single-system axioms.
\end{itemize}
\end{remark}

\subsection{Branch Factorization}

A key question: given that the composition axiom holds, how do branches
of the composite system $M(A,B)$ relate to branches of components $M_A$
and $M_B$?

\begin{definition}[Branch Factorization Map]
\label{def:factorization_map}
For a composite branch $\beta_{AB} \in \Br(M(A,B))$, define the
projection maps:
\begin{align}
\pi_A(\beta_{AB}) &= \beta_A = \{e \in \beta_{AB} : e \subseteq V_A\} \\
\pi_B(\beta_{AB}) &= \beta_B = \{e \in \beta_{AB} : e \subseteq V_B\}
\end{align}
The factorization map is:
\begin{equation}
\Phi: \Br(M(A,B)) \to \Br(M_A) \times \Br(M_B), \quad
\beta_{AB} \mapsto (\pi_A(\beta_{AB}), \pi_B(\beta_{AB}))
\end{equation}
\end{definition}

This map decomposes a composite branch into its constituent components.
The crucial question: Is $\Phi$ a bijection? We prove that it is, under
the hypotheses of \Cref{ax:composition}.

\section{Main Theorem: Tensor Product Structure}
\label{sec:tensor}

We now prove that the composition axiom induces tensor product structure through the factorization map.

\begin{theorem}[Tensor Product Structure from Composition Axiom]
\label{thm:tensor_product_main}
Let $M_A$ and $M_B$ be multiway systems satisfying \Cref{ax:composition}
(disjoint vertex sets, independent rules). Let $\Hilb_A$, $\Hilb_B$,
and $\Hilb_{AB}$ be the branch spaces (\Cref{def:branch_basis}) of
$M_A$, $M_B$, and $M(A,B)$ respectively, equipped with the
orthonormal inner product $\braket{\beta}{\beta'} = \delta_{\beta,\beta'}$.
Then:
\begin{equation}
\Hilb_{AB} \cong \Hilb_A \otimes \Hilb_B
\end{equation}
as inner product spaces, via the isomorphism
$\ket{\beta_{AB}} \mapsto \ket{\beta_A} \otimes \ket{\beta_B}$ where
$(\beta_A, \beta_B) = \Phi(\beta_{AB})$.
\end{theorem}

\begin{remark}[Hypotheses Made Explicit]
\Cref{thm:tensor_product_main} relies on three explicit hypotheses:
\begin{enumerate}
\item \textbf{Composition axiom} (\Cref{ax:composition}): Disjoint
  vertex sets and independent rules. This is an assumption, not derived
  from causal invariance.
\item \textbf{Orthonormal inner product}: The choice
  $\braket{\beta}{\beta'} = \delta_{\beta,\beta'}$ is essential for
  inner product factorization. Alternative inner products (e.g., based
  on hyperedge overlap) do not yield the factorization property.
\item \textbf{Causal invariance}: Ensures that the multiway evolution is
  well-defined and that branches represent operationally meaningful
  configurations (observer-independent outcomes).
\end{enumerate}
The composition axiom plays a role structurally analogous to the
composite system postulates in Hardy \cite{hardy2001quantum}, Chiribella
et al.\ \cite{chiribella2011informational}, and Masanes and M\"uller
\cite{masanes2011general}. Our framework does not eliminate this
postulate but rather gives it a concrete computational realization.
\end{remark}

The proof proceeds through the following lemmas and theorems:

\subsection{Branch Factorization Bijection}

\begin{lemma}[Branch Factorization Bijection]
\label{lem:factorization_bijection}
For independent systems $M_A, M_B$ with $V_A \cap V_B = \emptyset$, the factorization map $\Phi: \Br(M(A,B)) \to \Br(M_A) \times \Br(M_B)$ is a bijection.
\end{lemma}

\begin{proof}
We prove injectivity and surjectivity separately.

\textbf{Injectivity}: Suppose $\Phi(\beta_{AB}) = \Phi(\beta'_{AB})$. Then:
\begin{align}
\pi_A(\beta_{AB}) &= \pi_A(\beta'_{AB}) \\
\pi_B(\beta_{AB}) &= \pi_B(\beta'_{AB})
\end{align}

Since $V_A \cap V_B = \emptyset$, any hyperedge $e$ in $\beta_{AB}$ belongs to exactly one of three disjoint categories:
\begin{enumerate}
\item $e \subseteq V_A$ (edges involving only $A$ vertices),
\item $e \subseteq V_B$ (edges involving only $B$ vertices),
\item $e$ involves vertices from both $V_A$ and $V_B$.
\end{enumerate}

However, since rules in $\mathcal{R}_A$ only create hyperedges among $V_A$ vertices and rules in $\mathcal{R}_B$ only create hyperedges among $V_B$ vertices, category (3) is impossible. Thus:
\begin{equation}
\beta_{AB} = \pi_A(\beta_{AB}) \cup \pi_B(\beta_{AB})
\end{equation}

Therefore:
\begin{equation}
\beta_{AB} = \pi_A(\beta_{AB}) \cup \pi_B(\beta_{AB}) = \pi_A(\beta'_{AB}) \cup \pi_B(\beta'_{AB}) = \beta'_{AB}
\end{equation}
proving injectivity.

\textbf{Surjectivity}: Let $(\beta_A, \beta_B) \in \Br(M_A) \times \Br(M_B)$ be arbitrary. We must show there exists $\beta_{AB} \in \Br(M(A,B))$ with $\Phi(\beta_{AB}) = (\beta_A, \beta_B)$.

Since $\beta_A \in \Br(M_A)$, there exists a rule application sequence in $M_A$ producing configuration $\beta_A$. Similarly for $\beta_B$. Define:
\begin{equation}
\beta_{AB} = \beta_A \cup \beta_B
\end{equation}

This configuration is reachable in $M(A,B)$ by interleaving the rule sequences for $\beta_A$ and $\beta_B$ (both sequences commute since they operate on disjoint vertex sets). Thus $\beta_{AB} \in \Br(M(A,B))$, and:
\begin{align}
\pi_A(\beta_{AB}) &= \{e \in \beta_A \cup \beta_B : e \subseteq V_A\} = \beta_A \\
\pi_B(\beta_{AB}) &= \{e \in \beta_A \cup \beta_B : e \subseteq V_B\} = \beta_B
\end{align}
proving surjectivity.

Therefore $\Phi$ is a bijection.
\end{proof}

\subsection{Inner Product Factorization}

\begin{theorem}[Inner Product Factorization]
\label{thm:inner_product_factorization}
For the orthonormal branch bases defined in \Cref{def:branch_basis}, the composite inner product factorizes:
\begin{equation}
\braket{\beta_{AB}}{\beta'_{AB}} = \braket{\beta_A}{\beta'_A} \cdot \braket{\beta_B}{\beta'_B}
\end{equation}
where $(\beta_A, \beta_B) = \Phi(\beta_{AB})$ and $(\beta'_A, \beta'_B) = \Phi(\beta'_{AB})$.
\end{theorem}

\begin{proof}
By \Cref{def:branch_basis}, $\braket{\beta_{AB}}{\beta'_{AB}} = \delta_{\beta_{AB}, \beta'_{AB}}$ and similarly for $A$ and $B$ subsystems.

By \Cref{lem:factorization_bijection}, $\Phi$ is a bijection, so:
\begin{equation}
\beta_{AB} = \beta'_{AB} \iff (\beta_A, \beta_B) = (\beta'_A, \beta'_B) \iff (\beta_A = \beta'_A) \land (\beta_B = \beta'_B)
\end{equation}

Therefore:
\begin{align}
\delta_{\beta_{AB}, \beta'_{AB}} &=
\begin{cases}
1 & \text{if } (\beta_A = \beta'_A) \land (\beta_B = \beta'_B) \\
0 & \text{otherwise}
\end{cases} \\
&= \delta_{\beta_A, \beta'_A} \cdot \delta_{\beta_B, \beta'_B}
\end{align}

This completes the proof.
\end{proof}

\subsection{Tensor Product Isomorphism}

\begin{theorem}[Tensor Product Structure]
\label{thm:tensor_product}
For independent multiway systems $M_A, M_B$ with $V_A \cap V_B = \emptyset$, the composite branch space is isomorphic to the tensor product:
\begin{equation}
\Hilb_{AB} \cong \Hilb_A \otimes \Hilb_B
\end{equation}
\end{theorem}

\begin{proof}
Define the map $\psi: \Hilb_{AB} \to \Hilb_A \otimes \Hilb_B$ on basis elements by:
\begin{equation}
\psi(\ket{\beta_{AB}}) = \ket{\beta_A} \otimes \ket{\beta_B}
\end{equation}
where $(\beta_A, \beta_B) = \Phi(\beta_{AB})$, and extend linearly to all of $\Hilb_{AB}$.

\textbf{Well-defined}: By \Cref{lem:factorization_bijection}, $\Phi$ is a bijection, so $(\beta_A, \beta_B)$ is uniquely determined.

\textbf{Linear}: By construction.

\textbf{Bijective}: The basis $\{\ket{\beta_{AB}}\}$ is in bijection with $\Br(M_A) \times \Br(M_B)$ via $\Phi$, which corresponds to the basis $\{\ket{\beta_A} \otimes \ket{\beta_B}\}$ of $\Hilb_A \otimes \Hilb_B$.

\textbf{Inner product preserving}: For any $\ket{\gamma} = \sum_i c_i \ket{\beta_i}$ and $\ket{\eta} = \sum_j d_j \ket{\beta'_j}$:
\begin{align}
\braket{\psi(\gamma)}{\psi(\eta)} &= \left\langle \sum_i c_i (\ket{\beta_{A,i}} \otimes \ket{\beta_{B,i}}) \middle| \sum_j d_j (\ket{\beta'_{A,j}} \otimes \ket{\beta'_{B,j}}) \right\rangle \\
&= \sum_{i,j} c_i^* d_j \braket{\beta_{A,i}}{\beta'_{A,j}} \braket{\beta_{B,i}}{\beta'_{B,j}} \\
&= \sum_{i,j} c_i^* d_j \braket{\beta_i}{\beta'_j} \quad \text{(by \Cref{thm:inner_product_factorization})} \\
&= \braket{\gamma}{\eta}
\end{align}

Therefore $\psi$ is a unitary isomorphism, establishing $\Hilb_{AB} \cong \Hilb_A \otimes \Hilb_B$.
\end{proof}

\subsection{Computational Verification}

To validate \Cref{thm:tensor_product} empirically, we implemented multiway evolution for small hypergraph systems and tested the bijection $\Phi$ computationally.

\subsubsection{Test System}

We constructed two independent multiway systems:
\begin{itemize}
\item \textbf{System $A$}: Initial hypergraph $\{\{0,1\}\}$, rule $\{\{x,y\}\} \to \{\{x,y\}, \{y,z\}\}$ (vertex range $0$--$4$).
\item \textbf{System $B$}: Initial hypergraph $\{\{5,6\}\}$, rule $\{\{x,y\}\} \to \{\{x,y\}, \{y,z\}\}$ (vertex range $5$--$9$).
\end{itemize}

Both systems were evolved to depth $t=3$, generating:
\begin{itemize}
\item $\abs{\Br(M_A)} = 17$ branches for system $A$,
\item $\abs{\Br(M_B)} = 13$ branches for system $B$,
\item $\abs{\Br(M(A,B))} = 221$ composite branches.
\end{itemize}

\subsubsection{Bijection Tests}

\textbf{Injectivity Test}: For all $221$ composite branches $\beta_{AB}$, we computed $\Phi(\beta_{AB}) = (\beta_A, \beta_B)$ and verified uniqueness: no two distinct composite branches mapped to the same pair. Result: $100\%$ pass rate.

\textbf{Surjectivity Test}: We verified that $\abs{\{\Phi(\beta_{AB}) : \beta_{AB} \in \Br(M(A,B))\}} = \abs{\Br(M_A)} \times \abs{\Br(M_B)} = 17 \times 13 = 221$, confirming complete coverage. Result: $100\%$ pass rate.

\textbf{Inner Product Factorization Test}: We sampled $100$ random pairs $(\beta_{AB}, \beta'_{AB})$ and verified:
\begin{equation}
\abs{\braket{\beta_{AB}}{\beta'_{AB}} - \braket{\beta_A}{\beta'_A} \cdot \braket{\beta_B}{\beta'_B}} < 10^{-12}
\end{equation}
Result: $100\%$ pass rate, maximum error $< 10^{-15}$ (machine precision).

\subsubsection{Interpretation}

These computational results provide empirical support for \Cref{lem:factorization_bijection,thm:inner_product_factorization,thm:tensor_product}. While computational verification cannot substitute for formal proof (which we provide), it increases confidence and guards against implementation errors. The perfect agreement suggests the theoretical framework correctly captures the operational behavior of multiway systems.

\section{Local Discriminability Derivation}
\label{sec:ld}

Having established tensor product structure from operational composition, we now derive Chiribella's local discriminability axiom as a consequence.

\subsection{Chiribella's Local Discriminability Axiom}

We first state Chiribella et al.'s axiom precisely \cite{chiribella2011informational}:

\begin{axiom}[Local Discriminability (Chiribella)]
\label{ax:ld_chiribella}
Let $\rho_{AB}, \sigma_{AB}$ be states of a composite system $AB$ with the same marginal on subsystem $A$:
\begin{equation}
\Tr_B(\rho_{AB}) = \Tr_B(\sigma_{AB})
\end{equation}
Let $\ket{\psi_C}, \ket{\phi_C}$ be perfectly distinguishable pure states on an ancilla system $C$. Then the composite states:
\begin{equation}
\rho_{ABC} = \rho_{AB} \otimes \ketbra{\psi_C}{\psi_C}, \quad \sigma_{ABC} = \sigma_{AB} \otimes \ketbra{\phi_C}{\phi_C}
\end{equation}
are perfectly distinguishable.
\end{axiom}

In Chiribella's framework, this is an independent axiom. We now show it is a theorem given tensor product structure.

\subsection{Derivation from Tensor Product Structure}

\begin{theorem}[Local Discriminability from Tensor Product]
\label{thm:ld_from_tensor}
Assume the tensor product structure $\Hilb_{ABC} \cong \Hilb_{AB} \otimes \Hilb_C$ and perfect distinguishability of basis states. Then \Cref{ax:ld_chiribella} holds as a theorem.
\end{theorem}

\begin{proof}
Given $\ket{\psi_C}, \ket{\phi_C}$ perfectly distinguishable, we have $\braket{\psi_C}{\phi_C} = 0$ (orthogonality is equivalent to perfect distinguishability for pure states).

Define the measurement operator:
\begin{equation}
M = I_{AB} \otimes \ketbra{\psi_C}{\psi_C}
\end{equation}

The outcome probability for measuring $M$ on $\rho_{ABC}$ is:
\begin{align}
p_\psi(\rho) &= \Tr(M \rho_{ABC}) \\
&= \Tr[(I_{AB} \otimes \ketbra{\psi_C}{\psi_C})(\rho_{AB} \otimes \ketbra{\psi_C}{\psi_C})] \\
&= \Tr_{AB}(\rho_{AB}) \cdot \Tr_C(\ketbra{\psi_C}{\psi_C} \cdot \ketbra{\psi_C}{\psi_C}) \\
&= 1 \cdot \braket{\psi_C}{\psi_C}^2 = 1
\end{align}

Similarly, for $\sigma_{ABC}$:
\begin{align}
p_\psi(\sigma) &= \Tr(M \sigma_{ABC}) \\
&= \Tr[(I_{AB} \otimes \ketbra{\psi_C}{\psi_C})(\sigma_{AB} \otimes \ketbra{\phi_C}{\phi_C})] \\
&= \Tr_{AB}(\sigma_{AB}) \cdot \Tr_C(\ketbra{\psi_C}{\psi_C} \cdot \ketbra{\phi_C}{\phi_C}) \\
&= 1 \cdot \abs{\braket{\psi_C}{\phi_C}}^2 = 0
\end{align}

Therefore, the measurement $\{M, I - M\}$ gives outcome $\psi$ with probability $1$ for $\rho_{ABC}$ and probability $0$ for $\sigma_{ABC}$, perfectly distinguishing them.
\end{proof}

\begin{remark}
The key insight is that tensor product structure implies \textit{measurement factorization}: measurements on subsystem $C$ can be performed independently of $AB$, and their outcomes factorize via the tensor product inner product. This factorization property immediately yields local discriminability.
\end{remark}

\subsection{Non-Circularity Check}

A critical question: Does this derivation assume what it claims to prove?

\textbf{Check}:
\begin{enumerate}
\item Did we assume LD to prove tensor product structure (\Cref{thm:tensor_product})? \textbf{No}. The proof uses the composition axiom (\Cref{ax:composition}) and the factorization bijection (\Cref{lem:factorization_bijection}), neither of which reference LD. However, the composition axiom itself is an independent postulate, not derived from CI.
\item Did we assume tensor products before deriving them? \textbf{No}. Operational composition is defined as a union operation (vertices, hyperedges, rules), with no reference to tensor products.
\item Did we assume measurement factorization? \textbf{No}. Factorization follows from the definition of tensor product inner products, proven in \Cref{thm:inner_product_factorization}.
\end{enumerate}

Therefore, the derivation is non-circular: LD is proven from tensor product structure, which is proven from operational composition.

\subsection{Extension to Mixed States}

\Cref{thm:ld_from_tensor} applies to general (mixed) states $\rho_{AB}, \sigma_{AB}$. To see this, expand in the eigenbasis:
\begin{equation}
\rho_{AB} = \sum_k p_k \ketbra{\psi_k}{\psi_k}, \quad \sigma_{AB} = \sum_k q_k \ketbra{\phi_k}{\phi_k}
\end{equation}

Then:
\begin{align}
\rho_{ABC} &= \sum_k p_k \ketbra{\psi_k}{\psi_k} \otimes \ketbra{\psi_C}{\psi_C} \\
\sigma_{ABC} &= \sum_k q_k \ketbra{\phi_k}{\phi_k} \otimes \ketbra{\phi_C}{\phi_C}
\end{align}

The measurement $M = I_{AB} \otimes \ketbra{\psi_C}{\psi_C}$ yields:
\begin{align}
p_\psi(\rho) &= \sum_k p_k \Tr(\ketbra{\psi_k}{\psi_k}) \cdot \abs{\braket{\psi_C}{\psi_C}}^2 = \sum_k p_k = 1 \\
p_\psi(\sigma) &= \sum_k q_k \Tr(\ketbra{\phi_k}{\phi_k}) \cdot \abs{\braket{\psi_C}{\phi_C}}^2 = 0
\end{align}

Thus LD holds for mixed states as well.

\subsection{No-Signaling as Automatic Consequence}

The tensor product structure also implies \textit{no-signaling}: measurements on subsystem $A$ cannot affect outcome probabilities on subsystem $B$. This follows from factorization:
\begin{equation}
p_A(a \mid \rho_{AB}) = \Tr[(M_A^a \otimes I_B)\, \rho_{AB}] = \Tr_A[M_A^a\, \rho_A]
\end{equation}
where $\rho_A = \Tr_B(\rho_{AB})$ is the $A$-marginal.
The $A$-outcome probabilities depend only on $\rho_A$, not on any measurement performed on $B$, ensuring no-signaling. Therefore, no-signaling is built into tensor product structure, not an additional axiom.

\section{Discussion}
\label{sec:discussion}

\subsection{Honest Assessment of Contributions}

We begin with a frank assessment of what this work achieves and what it
does not.

\textbf{What is genuinely new.}
\begin{enumerate}
\item \textbf{Wolfram--operational QM bridge}: We construct an explicit
  mapping between multiway hypergraph branch spaces and the operational
  quantum mechanics framework of Chiribella et al.\
  \cite{chiribella2011informational}. This connection has been
  informally suggested \cite{wolfram2020new,gorard2020some} but not
  previously formalized with explicit proofs.
\item \textbf{Concrete tensor product construction}: Given the
  composition axiom (\Cref{ax:composition}), the branch factorization
  bijection (\Cref{lem:factorization_bijection}) and resulting tensor
  product isomorphism (\Cref{thm:tensor_product}) provide a clean,
  constructive proof within a specific computational model.
\item \textbf{LD as derived}: Within our framework, local
  discriminability (\Cref{thm:ld_from_tensor}) follows from tensor
  product structure, reducing Chiribella's axiom count by one.
\item \textbf{Purification test}: Computational validation that
  multiway branch spaces satisfy Chiribella's purification postulate.
\end{enumerate}

\textbf{What is not new.}
\begin{enumerate}
\item \textbf{Composition $\to$ tensor product}: The passage from
  independent subsystems to tensor product structure is the standard
  route in all operational QM frameworks. Chiribella et al.\
  \cite{chiribella2011informational} derive tensor products from their
  axioms including LD; Hardy \cite{hardy2001quantum} builds composite
  structure into his axiom set; Masanes and M\"uller
  \cite{masanes2011general} assume state spaces for composite systems.
  Our composition axiom plays an analogous structural role: it
  encodes the physical content (independent evolution) from which
  tensor products follow by standard mathematical arguments.
\item \textbf{LD from tensor products}: Once tensor product structure is
  established, the derivation of local discriminability
  (\Cref{thm:ld_from_tensor}) is straightforward and follows from
  standard properties of tensor product inner products. The conceptual
  value lies in recognizing LD as derivative, not in the proof
  technique.
\item \textbf{Orthonormality structure}: The inner product
  $\braket{\beta}{\beta'} = \delta_{\beta,\beta'}$ does much of the
  mathematical work. This is a strong assumption, and the factorization
  property relies on it essentially.
\end{enumerate}

\subsection{Comparison with Existing Frameworks}

\begin{table}[htbp]
\centering
\caption{Axiom comparison across operational QM frameworks. All
  frameworks assume some form of composition for composite systems.}
\label{tab:axiom_comparison}
\begin{tabular}{@{}lcccc@{}}
\toprule
\textbf{Framework} & \textbf{Axioms} & \textbf{LD Status} &
  \textbf{Composition} & \textbf{Foundation} \\
\midrule
Chiribella et al.\ \cite{chiribella2011informational} & 6 & Axiom &
  Assumed & Operational \\
Hardy \cite{hardy2001quantum} & 5 & Implicit & Simplicity axiom &
  Operational \\
Masanes \& M\"uller \cite{masanes2011general} & 5 & Info.\ locality &
  State space axiom & Operational \\
This work & \textbf{5} & \textbf{Theorem} & \textbf{Composition axiom} &
  CI + composition \\
\bottomrule
\end{tabular}
\end{table}

\Cref{tab:axiom_comparison} makes explicit a point that deserves
emphasis: \textit{every} operational QM reconstruction assumes
composition in some form. Our framework does not eliminate this
assumption; rather, it gives composition a concrete computational
realization (disjoint multiway systems) and connects it to discrete
causal structures.

The meaningful differences are:
\begin{itemize}
\item \textbf{LD derived}: We reduce Chiribella's six axioms to five by
  deriving LD from tensor product structure. This is a genuine (if
  modest) simplification.
\item \textbf{Computational foundation}: Unlike purely axiomatic
  frameworks, our multiway systems provide an explicit computational
  model that can be simulated, tested, and extended.
\item \textbf{Causal invariance connection}: Our framework connects to
  the Wolfram Physics Project's program of deriving physics from
  discrete causal structures, providing a formal link between that
  program and established operational quantum mechanics.
\end{itemize}

\subsection{The Categorical Perspective and Its Gaps}

Gorard, Namuduri, and Arsiwalla \cite{gorard2021categorical} have shown that
the multiway evolution graph of a causally invariant system, viewed
categorically, forms a symmetric monoidal category (SMC). This connects
to Meseguer's \cite{meseguer1992conditional} result that concurrent rewriting with confluence
gives rise to SMC structure.

However, the categorical path from causally invariant multiway systems
to the category of finite-dimensional Hilbert spaces ($\mathbf{FdHilb}$)
requires three steps:
\begin{enumerate}
\item CI $\to$ SMC: established by Gorard et al.\ and Meseguer.
\item SMC $\to$ dagger compact category: requires a dagger functor
  (encoding time-reversal/adjoint structure) and compact closure
  (encoding entanglement via cup/cap morphisms). These structures are
  \textit{not} supplied by causal invariance.
\item Dagger compact category $\to$ $\mathbf{FdHilb}$: established by
  the Abramsky--Coecke \cite{abramsky2004categorical} representation theorem.
\end{enumerate}

The gap at step 2 is where specifically quantum features enter. Our
composition axiom addresses a related but distinct gap: it provides
tensor product structure on \textit{state spaces} (not on the category
of processes), enabling the derivation of LD. The categorical program
and our state-space approach are complementary; a complete derivation of
quantum mechanics from causal invariance would require resolving both.

\subsection{Novelty Assessment}

We estimate our contribution's novelty at approximately $15$--$20\%$:

\begin{itemize}
\item \textbf{Wolfram--operational QM bridge ($\sim 10$--$12\%$)}:
  Connecting multiway systems to Chiribella's framework via explicit
  mappings (states $\leftrightarrow$ branches, composition
  $\leftrightarrow$ disjoint union, measurements $\leftrightarrow$
  projections). This bridge has been informally discussed but not
  previously formalized.
\item \textbf{LD derivation ($\sim 3$--$5\%$)}: Recognizing and proving
  that LD follows from tensor product structure. Mathematically
  straightforward, but conceptually clarifying.
\item \textbf{Computational verification ($\sim 2$--$3\%$)}: Empirical
  testing adds reproducibility and guards against errors.
\end{itemize}

The composition-to-tensor-product passage itself ($0\%$ novelty) is
standard in operational QM. Our framework instantiates it in a new
setting but does not introduce a new method.

This honest assessment reflects the fact that most components build on
existing ideas. The value lies in the synthesis: bringing together two
research programs (Wolfram physics and operational QM) that have
developed independently, and demonstrating their compatibility within
a concrete computational framework.

\subsection{What Causal Invariance Does and Does Not Provide}

To avoid misunderstanding, we state explicitly the role of causal
invariance in our framework:

\textbf{CI provides:}
\begin{itemize}
\item Observer-independent causal structure (the causal partial order is
  gauge-invariant).
\item Well-defined multiway evolution (confluence ensures meaningful
  branch spaces).
\item A foundation for operationally meaningful states (branches as
  maximally refined outcomes).
\end{itemize}

\textbf{CI does not provide:}
\begin{itemize}
\item Subsystem decomposition (how to carve one system into independent
  parts).
\item Disjoint composition (the existence of non-interacting regions).
\item Tensor product structure (which requires composition as an
  additional input).
\item Complex Hilbert space structure (our framework is real-valued).
\end{itemize}

This separation is not a deficiency unique to our approach. In general
relativity, the Einstein equations determine local dynamics but not
global topology; the manifold structure is an independent input.
Similarly, in our framework, causal invariance determines evolution
dynamics but not compositional structure.

\section{Limitations and Future Work}
\label{sec:limitations}

\subsection{Causal Invariance Does Not Force Disjoint Composition}
\label{sec:lim_composition}

The most significant limitation of this work is that our central
structural assumption---the composition axiom (\Cref{ax:composition})---is
not derived from causal invariance. This independence has been established
through three converging lines of evidence:

\begin{enumerate}
\item \textbf{Counterexamples}: Causally invariant hypergraph rewriting
  systems can be fully connected, with every vertex participating in
  rules that span the entire graph. Such systems admit no partition
  $V = V_A \sqcup V_B$ for which the rules decompose into independent
  subsets. Thus causal invariance alone does not guarantee the existence
  of independently composable subsystems.

\item \textbf{Categorical analysis}: While Gorard et al.\ \cite{gorard2021categorical} show that
  causally invariant multiway systems form a symmetric monoidal category
  (SMC) via the Meseguer \cite{meseguer1992conditional} construction, the monoidal product in this
  category acts on \textit{processes} (rewrite sequences), not on
  \textit{state spaces}. The passage from SMC to dagger compact
  categories---and thence to $\mathbf{FdHilb}$ via the Abramsky--Coecke
  \cite{abramsky2004categorical} representation theorem---requires a dagger functor and compact closure,
  neither of which is supplied by causal invariance.

\item \textbf{Literature}: No published proof derives tensor product
  structure on state spaces from causal invariance alone. Zanardi, Lidar,
  and Lloyd \cite{zanardi2004quantum} have shown that tensor product structures are
  observable-induced (relative to accessible measurements), further
  suggesting that compositional structure requires input beyond dynamics.
\end{enumerate}

This limitation is shared, in structural form, with all operational QM
reconstructions: composition must be assumed. Our framework makes this
assumption explicit rather than implicit, which we regard as a virtue of
the formulation.

\subsection{Composition Axiom as Additional Postulate}

The composition axiom (\Cref{ax:composition}) is analogous to the
following postulates in other contexts:

\begin{itemize}
\item \textbf{Manifold structure in GR}: The Einstein equations determine
  local gravitational dynamics but not the topology of the spacetime
  manifold. Similarly, CI determines evolution dynamics but not
  compositional structure.
\item \textbf{Split property in AQFT}: In algebraic quantum field theory,
  the split property for spacelike-separated von Neumann algebras
  guarantees tensor product decomposition. This property is an additional
  condition, not derived from the field equations alone.
\item \textbf{Composite system axiom in operational QM}: Every
  reconstruction of quantum theory from operational axioms
  \cite{hardy2001quantum,chiribella2011informational,masanes2011general}
  assumes some form of compositional structure for joint systems.
\end{itemize}

A promising direction for future work is investigating whether
composition might emerge in a suitable limit. In the Wolfram model, if
large causally invariant hypergraphs generically develop emergent spatial
locality (as conjectured), then spacelike-separated regions would
naturally satisfy the composition axiom. However, each step in this
argument---emergent locality, sufficient spatial separation,
persistence of independence---remains an open problem.

\subsection{Real-Valued Framework}

Our construction yields tensor products of \textit{real} inner product
spaces. Chiribella et al.'s full axiom set, combined with local
discriminability, selects \textit{complex} quantum mechanics
specifically, ruling out real and quaternionic alternatives. Our
framework does not address this selection:

\begin{itemize}
\item The branch space inner product
  $\braket{\beta}{\beta'} = \delta_{\beta,\beta'}$ is real-valued.
\item No mechanism for complex phases arises from our axioms.
\item The distinction between real, complex, and quaternionic quantum
  mechanics requires additional structure (e.g., Chiribella's local
  discriminability + purification jointly select $\mathbb{C}$).
\end{itemize}

\textbf{Future work}: Investigate whether complex phases emerge from
additional properties of multiway evolution (e.g., from the branchial
graph structure or from causal cone interference patterns).

\subsection{Inner Product Structure}

The orthonormality assumption
$\braket{\beta}{\beta'} = \delta_{\beta,\beta'}$ is essential to our
results. The inner product factorization
(\Cref{thm:inner_product_factorization}) relies on the Kronecker delta
structure: the product of Kronecker deltas factors as a single Kronecker
delta if and only if the bijection $\Phi$ maps distinct composite
branches to distinct pairs.

Alternative inner products---for instance, based on hyperedge overlap
between configurations---do not yield this factorization property. We
tested an edge-overlap inner product
$\braket{\beta}{\beta'}_{\text{edge}} = |E(\beta) \cap E(\beta')| /
|E(\beta) \cup E(\beta')|$, which failed to factorize and did not
support tensor product structure. The orthonormal inner product is thus
a substantive choice, not merely a convenient one.

\textbf{Future work}: Explore whether the orthonormal inner product
can be derived from operational considerations (e.g., as the unique
inner product compatible with perfect distinguishability and some
continuity condition).

\subsection{Restriction to Independent Systems}

Our results apply to systems with disjoint vertex sets
($V_A \cap V_B = \emptyset$), representing \textit{independent}
evolution. This excludes:
\begin{itemize}
\item \textbf{Entangled states}: States not factorizable as
  $\ket{\beta_A} \otimes \ket{\beta_B}$ require shared causal structure.
  Extending our framework to entanglement is an important open problem.
\item \textbf{Interacting systems}: Rules coupling $A$ and $B$ vertices
  would violate the composition axiom. Modeling interactions requires
  relaxing \Cref{ax:composition}.
\end{itemize}

\textbf{Future work}: Investigate whether controlled interactions
(partial vertex overlap, or rules that occasionally span both regions)
give rise to entanglement in a natural way.

\subsection{Continuum Limit}

We work with discrete multiway systems, not continuous Hilbert spaces.
The connection to standard quantum mechanics requires taking a continuum
limit:
\begin{equation}
\text{Discrete multiway system} \xrightarrow{\text{limit}}
\text{Continuous QM}
\end{equation}

This limit is \textit{assumed}, not derived. A systematic scaling study
of $13$ Wolfram rules across diverse types~\cite{companion_lovelock}
found that all dynamically expanding rules exhibit $\kappa \propto 1/N
\to 0$, empirically disconfirming the continuum limit for generic
causal-invariant systems. Only three small static rules ($N \leq 8$)
showed nonzero convergence. Until this barrier is resolved, the
discrete-to-continuous passage remains a significant open problem.

Our results hold at the discrete level: tensor product structure and
local discriminability are properties of the multiway branch spaces
themselves and do not require a continuum limit.

\textbf{Future work}: Investigate whether restricted classes of
causal-invariant systems (e.g., those with specific growth properties)
admit a continuum limit, possibly using techniques from lattice QFT or
noncommutative geometry.

\subsection{Computational vs.\ Formal Proofs}

Our tensor product isomorphism (\Cref{thm:tensor_product}) is proven
mathematically, with computational verification supplementing the proof
by testing specific instances. While the computational results ($100\%$
accuracy over $221$ states) increase confidence, they do not constitute
a formal proof for all multiway systems.

\textbf{Future work}: Formalize the proofs in a proof assistant (e.g.,
Lean 4) to provide machine-verified certainty.

\subsection{Scope of Derivation}

We derive LD within the \textit{Hilbert space} framework, not the more
general categorical frameworks used by Coecke et al.\
\cite{coecke2010quantum}. Our results may not apply to generalized
probabilistic theories or non-Hilbert operational theories.

\textbf{Future work}: Extend the derivation to categorical quantum
mechanics, investigating whether multiway composition realizes the
categorical tensor in dagger compact categories.

\section{Conclusion}
\label{sec:conclusion}

We have demonstrated that the tensor product structure of composite quantum systems and the principle of local discriminability emerge as mathematical consequences of two independent principles: causal invariance and a composition axiom for independent subsystems. Our main results are:

\begin{enumerate}
\item \textbf{Tensor product from composition axiom} (\Cref{thm:tensor_product}): For independent systems satisfying the composition axiom (disjoint vertex sets, independent rules), the composite branch space is isomorphic to the tensor product of component spaces through an explicit factorization bijection (\Cref{lem:factorization_bijection}).

\item \textbf{LD as derived theorem} (\Cref{thm:ld_from_tensor}): Local discriminability follows from tensor product structure combined with perfect distinguishability, reducing Chiribella et al.'s six axioms to five.

\item \textbf{Computational verification}: Empirical testing on $221$ composite states confirms the bijection and inner product factorization with $100\%$ accuracy, supporting the theoretical framework.

\item \textbf{Explicit composition axiom}: Operational composition is introduced as a clearly stated axiom, making the logical dependencies between causal invariance and compositional structure transparent. This axiom is physically motivated by spacelike separation but not derived from causal invariance alone.
\end{enumerate}

These results provide an explicit bridge between Wolfram multiway systems and operational quantum mechanics, demonstrating that multiway branch spaces equipped with the composition axiom satisfy the operational axioms of quantum theory.

\subsection{Honest Assessment}

In the spirit of intellectual honesty, we acknowledge:
\begin{itemize}
\item \textbf{Novelty ($\sim 15$--$20\%$)}: Most components build on existing ideas; novelty lies in synthesis and explicit formalization of the Wolfram--operational QM connection.
\item \textbf{Composition axiom}: The tensor product structure requires an independent composition axiom, not derived from causal invariance. This limitation is shared with all operational QM frameworks.
\item \textbf{Limitations}: Real-valued framework (not complex), restriction to independent systems, discrete framework only (continuum limit is an open problem with mixed evidence), orthonormal inner product essential.
\item \textbf{Computational proofs}: Verification complements but does not replace formal proof; independent replication needed.
\end{itemize}

\subsection{Future Directions}

The most promising directions for extending this work include:
\begin{enumerate}
\item \textbf{Entanglement}: Generalizing to systems with shared causal structure.
\item \textbf{Continuum limit}: Identifying classes of causal-invariant systems that admit a continuum limit to standard QM, given that generic expanding rules exhibit vanishing convergence.
\item \textbf{Other axioms}: Deriving purification, ideal compression, or other operational axioms from causal invariance.
\item \textbf{Experimental tests}: Identifying empirical signatures distinguishing this approach from standard QM.
\item \textbf{Gravitational coupling}: Connecting to Wolfram's Einstein equations derivation to unify QM and GR.
\end{enumerate}

\subsection{Final Remarks}

By deriving quantum theory's composite system structure from two independent operational principles---causal invariance and disjoint composition---we take a step toward understanding why quantum mechanics has the form it does. The explicit separation of these principles clarifies what causal invariance alone provides and what requires additional structure. Whether this approach ultimately contributes to fully reconstructing quantum theory remains an open question requiring continued investigation, computational testing, and independent verification.

We invite the community to scrutinize our claims, replicate our computational results, and explore the implications for quantum foundations and beyond.

\section*{Acknowledgments}

We thank the Wolfram Physics Project community for foundational ideas on causal invariance and multiway systems. Conversations with researchers in operational quantum mechanics have shaped our understanding of Chiribella's framework. All errors and limitations are the author's responsibility.

\section*{Data Availability}

All computational code, test data, and verification scripts are available at:
\begin{center}
\url{https://github.com/MaxZhuravlev/physics-operational-qm-ci}
\end{center}

\begin{thebibliography}{99}

\bibitem{chiribella2011informational}
G. Chiribella, G. M. D'Ariano, and P. Perinotti,
\textit{Informational derivation of quantum theory},
Phys. Rev. A \textbf{84}, 012311 (2011).
arXiv:1011.6451 [quant-ph].

\bibitem{hardy2001quantum}
L. Hardy,
\textit{Quantum theory from five reasonable axioms},
arXiv:quant-ph/0101012 (2001).

\bibitem{masanes2011general}
Ll. Masanes and M. P. Müller,
\textit{A derivation of quantum theory from physical requirements},
New J. Phys. \textbf{13}, 063001 (2011).
arXiv:1004.1483 [quant-ph].

\bibitem{dakic2009quantum}
B. B.\ Daki\'{c} and \v{C}.\ Brukner,
\textit{Quantum theory and beyond: is entanglement special?}
in \textit{Deep Beauty: Understanding the Quantum World through Mathematical Innovation},
H. Halvorson (ed.), Cambridge University Press, 2011.
arXiv:0911.0695 [quant-ph].

\bibitem{wolfram2020new}
S. Wolfram,
\textit{A New Kind of Science},
Wolfram Media, 2002;
\textit{A Project to Find the Fundamental Theory of Physics},
Wolfram Media, 2020.
\url{https://www.wolframphysics.org/}.

\bibitem{gorard2020some}
J. Gorard,
\textit{Some Relativistic and Gravitational Properties of the Wolfram Model},
arXiv:2004.14810 [gr-qc] (2020).

\bibitem{coecke2010quantum}
B. Coecke and A. Kissinger,
\textit{Picturing Quantum Processes: A First Course in Quantum Theory and Diagrammatic Reasoning},
Cambridge University Press, 2017.

\bibitem{gorard2021categorical}
J. Gorard, M. Namuduri, and X. D. Arsiwalla,
\textit{Fast Automated Reasoning over String Diagrams using Multiway Causal Structure},
arXiv:2105.04057 [cs.SC] (2021).

\bibitem{meseguer1992conditional}
J. Meseguer,
\textit{Conditional rewriting logic as a unified model of concurrency},
Theoretical Computer Science \textbf{96}, 73--155 (1992).

\bibitem{zanardi2004quantum}
P. Zanardi, D. A. Lidar, and S. Lloyd,
\textit{Quantum tensor product structures are observable induced},
Phys. Rev. Lett. \textbf{92}, 060402 (2004).
arXiv:quant-ph/0308043.

\bibitem{haag1992local}
R. Haag,
\textit{Local Quantum Physics: Fields, Particles, Algebras},
Springer-Verlag, Berlin, 1992.

\bibitem{abramsky2004categorical}
S. Abramsky and B. Coecke,
\textit{A categorical semantics of quantum protocols},
in \textit{Proceedings of the 19th Annual IEEE Symposium on Logic in Computer Science}, 415--425 (2004).
arXiv:quant-ph/0402130.

\bibitem{companion_lovelock}
M. Zhuravlev,
\textit{Where the Lovelock Bridge Breaks: Negative Results and New Directions for Connecting Discrete and Continuous Spacetime Emergence},
Preprint (2026).

\end{thebibliography}

\end{document}
